\documentclass[11pt,a4paper]{article}

\IfFileExists{glyphtounicode.tex}{\input{glyphtounicode}\pdfgentounicode=1}{}

\usepackage[utf8]{inputenc}
\usepackage[T1]{fontenc}
\IfFileExists{lmodern.sty}{\usepackage{lmodern}}{}
\usepackage[margin=25mm]{geometry}
\usepackage{amsmath}
\usepackage{amssymb}
\usepackage{booktabs}
\usepackage{longtable}
\usepackage{array}
\usepackage{xcolor}
\usepackage{hyperref}
\IfFileExists{microtype.sty}{\usepackage[protrusion=true,expansion=false]{microtype}}{}
\IfFileExists{enumitem.sty}{\usepackage{enumitem}}{}
\providecommand{\DisableLigatures}[1]{}
\AtBeginDocument{\DisableLigatures{encoding=*,family=*}}

\hypersetup{
  colorlinks=true,
  linkcolor=black,
  citecolor=black,
  urlcolor=blue,
  pdfauthor={Harun Aktas},
  pdftitle={DAD FieldWorks: An Evidence Contract Framework for Trustworthy Computational Electromagnetics},
  pdfsubject={Public technical whitepaper},
  pdfkeywords={computational electromagnetics, evidence contracts, RF, microwave, resonator, claim boundaries}
}

\setlength{\parindent}{0pt}
\setlength{\parskip}{0.62em}
\renewcommand{\arraystretch}{1.18}
\emergencystretch=2em

\newcommand{\code}[1]{\texttt{#1}}
\newcommand{\statusbox}[1]{%
  \vspace{0.55em}
  \noindent\fcolorbox{black}{gray!7}{%
    \begin{minipage}{0.94\linewidth}
    #1
    \end{minipage}}
  \vspace{0.55em}
}
\newcommand{\claimline}[1]{\textit{#1}}

\title{\textbf{DAD FieldWorks}\\
{\large An Evidence Contract Framework for Trustworthy Computational Electromagnetics}\\
\vspace{0.4em}
{\normalsize Validation-Aware RF, Microwave Resonator, and Future Physics-AI Solver Workflows}}
\author{Harun Aktas}
\date{Technical Whitepaper / Public Specification Draft\\Version 0.8 draft\\June 2026}

\begin{document}

\maketitle

\statusbox{\textbf{Document status:} Public architecture draft.\\
\textbf{Claim status:} No external validation claim. No production readiness claim. No commercial solver equivalence claim.\\
\textbf{Copyright:} Copyright \copyright{} 2026 Harun Aktas. All rights reserved.\\
\textbf{License notice:} No license is granted for reuse, redistribution, or derivative work without written permission.}

\statusbox{\textbf{Front-matter claim boundary.}
This document describes a private research and development project through public-safe architecture notes, bounded diagnostic examples, public visualizations, and roadmap material. It does not publish private solver source code, private validation workbench files, private project-control files, notebooks, copyrighted reference text, or implementation internals.}

\begin{abstract}
Computational electromagnetics (CEM) results can look plausible while still being wrong. A field plot, modal surface, spectrum, or time trace can be visually persuasive without proving residual quality, boundary consistency, reference agreement, reproducibility, or validation status. DAD FieldWorks is being developed as an evidence-contract framework that separates numerical output from what that output is allowed to claim. The framework attaches evidence records, residual checks, reference-comparison records, gate rows, decision classes, failure classes, risk rows, and claim boundaries to computational results. The current public state is bounded internal Alpha evidence, not a production solver and not an externally validated system. The longer-term direction is to make the same evidence discipline usable around classical numerical solvers, future discontinuous Galerkin time-domain (DGTD) solvers, and future physics-AI backends such as physics-informed neural networks (PINNs) or neural operators.
\end{abstract}

\textbf{Keywords:} computational electromagnetics; evidence contracts; RF; microwave resonators; PEC cavity; Yee incidence; residual checks; claim boundaries; physics-AI.

\tableofcontents
\newpage

\section{Introduction}

Computational electromagnetics produces numerical artifacts that can be useful long before they are safe to use as validation or production evidence. A field plot can look physically plausible, a spectrum can exhibit a strong peak, and an eigenvalue can be finite, yet the result may still depend on unresolved mapping assumptions, incomplete boundary checks, insufficient residual policy, or unreviewed comparison scope.

DAD FieldWorks addresses this as an architecture problem. The project is being shaped into an evidence-controlled computational electromagnetics framework for validation-aware RF, microwave resonator, cavity, and future physics-AI workflows. Its central premise is that a numerical result should not be trusted simply because it is computable or visually convincing. It should be trusted only to the extent that its evidence record supports a bounded claim.

The project is not presented here as a finished solver product. The current public material describes a private research and development project, public diagnostic visualizations, and a public specification draft. The current public-safe milestone is a bounded internal PEC cavity eigenmode evidence path with residual and analytical reference comparison. Stronger public claims remain closed.

\section{Why Field Plots Are Not Enough}

A field plot can be beautiful and still be wrong. A result can converge numerically and still violate boundary assumptions. A neural surrogate can be fast and still hallucinate plausible physical fields. In each case, the visual or numerical surface is only a candidate artifact.

For CEM workflows, the missing evidence may include:

\begin{itemize}
  \item residual quality,
  \item boundary-condition consistency,
  \item finite-value and energy sanity,
  \item component-location and grid-mapping policy,
  \item reference-problem equivalence,
  \item reproducibility records,
  \item decision and failure classification,
  \item claim authorization.
\end{itemize}

DAD FieldWorks therefore treats solver output, surrogate output, field slices, spectra, and visual diagnostics as untrusted until evidence gates permit a stronger interpretation. This is especially important for future physics-AI workflows, where plausible fields can be produced quickly but still require Maxwell residual checks, boundary checks, divergence checks, energy checks, and reference comparison when available.

\section{Evidence Contracts for Trustworthy Computational Electromagnetics}

An \emph{evidence contract} is a structured envelope around a numerical or diagnostic result. It couples the payload with explicit evidence state and claim state. The payload may be a frequency estimate, field slice, time trace, spectrum, residual, matrix diagnostic, or replay artifact. The contract records what the artifact is, what assumptions it uses, what gates it passed or failed, and what the artifact is allowed to claim.

The main concepts are:

\begin{description}
  \item[Evidence Contract] A result envelope that binds numerical output to evidence state, decision state, risk state, comparison state, and claim boundaries.
  \item[Evidence Row] A factual record about an artifact, such as geometry, grid, field component, residual value, finite-value check, or source/probe condition.
  \item[Gate Row] A pass, fail, blocked, or not-applicable condition that controls whether a stronger interpretation is allowed.
  \item[Decision Class] A compact status for the artifact, such as diagnostic, inconclusive, blocked, accepted within a bounded scope, or rejected.
  \item[Failure Class] A structured reason why a result failed or could not be promoted.
  \item[Risk Row] A record of unresolved ambiguity or interpretation risk.
  \item[Reference Comparison Row] A record tying internal output to a stated reference problem, reference value, equivalence policy, and comparison metric.
  \item[Claim Boundary] A public or internal statement describing what must not be inferred from the current evidence.
  \item[Required Future Work Row] A record of missing evidence needed before a stronger gate may open.
  \item[Artifact Replay] A way to re-read and re-check recorded evidence without re-running private simulation paths.
  \item[\code{ProductionAllowedQ}] A conservative flag that remains false unless a separate production gate supports changing it.
\end{description}

\code{ExternalValidationQ} remains false by default. \code{ProductionAllowedQ} remains false by default. Numerical success does not automatically promote a claim.

\section{Current Public-Safe Alpha State}

DAD FieldWorks is a private research and development project. The current public-safe state can be summarized as follows:

\begin{itemize}
  \item The current public-safe milestone is a bounded PEC cavity eigenmode evidence path.
  \item Native Yee curl incidence microprototype work may be described as bounded, internal, operator-near evidence.
  \item FDTD resonator ringdown and FFT spectrum assets are public diagnostic visualizations.
  \item Microwave Cavity Eigenmode Birth is a public deterministic visualization of an evidence path.
  \item The C++ direction is future product-core and artifact-replay direction, not a production electromagnetic solver claim.
  \item No graphical application is claimed.
  \item No external validation claim is made.
  \item No production readiness claim is made.
\end{itemize}

The public repository does not contain private solver source code, notebooks, private validation artifacts, or internal project-control files. The public materials are designed to communicate the evidence-contract idea without publishing private implementation details.

\section{Evidence Row and Gate Row Model}

The evidence row and gate row split is the practical center of the framework. Evidence rows record what is known. Gate rows decide what can be claimed from it.

\begin{center}
\begin{tabular}{p{0.26\linewidth}p{0.38\linewidth}p{0.26\linewidth}}
\toprule
Concept & Purpose & Example \\
\midrule
Evidence Row & Record factual artifact state & Grid size, residual norm, finite-value check \\
Gate Row & Control claim promotion & Residual gate passes or blocks \\
Reference Comparison Row & Bind result to reference problem & PEC cavity frequency comparison \\
Risk Row & Preserve unresolved ambiguity & Mapping policy remains material \\
Claim Boundary Row & State forbidden interpretations & Diagnostic only; no validation claim \\
Required Future Work Row & Identify missing evidence & Divergence policy review required \\
\bottomrule
\end{tabular}
\end{center}

This model keeps a finite number from becoming an unsupported engineering conclusion. A result can be useful, reproducible, and still bounded to diagnostic status.

\section{Claim Boundary Model}

Claim boundaries are first-class records. They prevent an internally useful result from being promoted to public validation, production readiness, or solver equivalence without evidence.

\begin{center}
\begin{tabular}{p{0.43\linewidth}p{0.43\linewidth}}
\toprule
Allowed current public claim & Forbidden claim \\
\midrule
Bounded internal prototype & External validation \\
Diagnostic visualization & Production readiness \\
Public roadmap & Commercial solver equivalence \\
Evidence contract architecture & Complete quantum hardware solver \\
Validation-aware research direction & Qubit simulation \\
Future DGTD route & Production DGTD backend \\
Future PINN route & Production AI solver \\
\bottomrule
\end{tabular}
\end{center}

The boundary is intentionally repetitive at decision points: the project may show meaningful research progress while still keeping validation and production gates closed.

\section{Classical Solver Path: PEC Cavity Evidence Chain}

The PEC cavity case study gives a controlled classical reference path for evidence-contract development. The public-safe rectangular cavity parameters are:

\begin{center}
\begin{tabular}{ll}
\toprule
Quantity & Value \\
\midrule
Cavity dimensions & $a=0.08\,\mathrm{m}$, $b=0.084\,\mathrm{m}$, $d=0.084\,\mathrm{m}$ \\
Mode indices & $m=n=p=1$ \\
Speed of light & $c_0=299792458\,\mathrm{m/s}$ \\
Analytical reference frequency & approximately $3.143165987503105\,\mathrm{GHz}$ \\
\bottomrule
\end{tabular}
\end{center}

For a rectangular cavity with simple PEC boundaries, the analytical frequency used as a reference can be written in the familiar separated form
\begin{equation}
f_{mnp} =
\frac{c_0}{2}
\sqrt{\left(\frac{m}{a}\right)^2+
      \left(\frac{n}{b}\right)^2+
      \left(\frac{p}{d}\right)^2}.
\label{eq:pec-reference-frequency}
\end{equation}

The scalar Helmholtz refinement diagnostic remains a numerical record, not a vector Maxwell validation claim:

\begin{center}
\begin{tabular}{lllll}
\toprule
Level & Grid & Frequency (GHz) & Relative error & Residual norm \\
\midrule
0 & $\{9,8,8\}$ & 3.128307359058 & 0.0047272809 & $2.63\times10^{-15}$ \\
1 & $\{18,16,16\}$ & 3.139012096545 & 0.0013215627 & $1.29\times10^{-14}$ \\
2 & $\{27,24,24\}$ & 3.141247276062 & 0.0006104391 & $1.68\times10^{-14}$ \\
\bottomrule
\end{tabular}
\end{center}

The same-grid scalar diagnostic remains claim-bounded because mapping policy ambiguity is material:

\begin{center}
\begin{tabular}{lllll}
\toprule
Policy & Matrix size & Frequency (GHz) & Relative error & Residual norm \\
\midrule
1 & $504\times504$ & 2.778104015200 & 0.1161446687 & $1.83\times10^{-15}$ \\
2 & $336\times336$ & 3.122528486251 & 0.0065658325 & $1.14\times10^{-15}$ \\
3 & $504\times504$ & 3.126946915965 & 0.0051601066 & $2.37\times10^{-15}$ \\
\bottomrule
\end{tabular}
\end{center}

The scalar path does not become a full vector Maxwell eigenmode claim. The same-grid ambiguity remains a claim boundary. Residual, mass, divergence or gauge, component-location, and incidence-compatibility policies must remain explicit before stronger claims are allowed.

\section{Native Yee Curl Incidence Microprototype}

The native Yee curl incidence microprototype is an operator-near public-safe diagnostic. It records Yee electric unknowns, local orientation, signed incidence rows, and bounded row construction. Its role is to clarify how local edge orientation can enter a curl incidence row.

A bounded internal audit path indicates that the microprototype is internally claim-safe within its limited scope. That scope is narrow:

\begin{itemize}
  \item it is not curl-curl assembly,
  \item it is not a production incidence matrix,
  \item it is not an eigensolve,
  \item it is not an FDTD rerun,
  \item it is not CPML,
  \item it is not production solver behavior.
\end{itemize}

The value of the microprototype is architectural: it makes operator evidence explicit before larger matrices or solver layers are allowed to inherit stronger claims.

\section{FDTD Resonator Ringdown and Spectral Diagnostics}

The public FDTD resonator visualization shows time-domain pulse excitation, field build-up, ringdown, a probe trace, and an FFT spectrum derived from the probe signal. It is a diagnostic example for resonator-oriented RF workflows.

The microwave resonator framing is relevant to quantum-hardware-oriented RF engineering because cavities, feedlines, resonators, and packaging modes are classical electromagnetic structures that can influence hardware behavior. This does not make the diagnostic a qubit simulation, Josephson junction model, processor-level quantum-hardware validation, external validation, or production readiness claim.

\section{Microwave Cavity Eigenmode Birth Visualization}

The Microwave Cavity Eigenmode Birth visualization is a deterministic public visualization of an evidence path. It connects Yee electric unknowns, oriented curl incidence, bounded curl-curl structure, a standing PEC cavity eigenmode slice, residual check, and analytical reference comparison.

Its purpose is explanatory. It shows how local orientation records and bounded diagnostic rows can lead toward a cavity-mode evidence chain. It does not show a production solver, does not publish private implementation details, and does not make an external validation claim.

\section{Future Backend Routes: Classical, DGTD, and Physics-AI}

The long-term architecture is backend-agnostic. A result may originate from a classical solver, a future DGTD route, or a future physics-AI route. In each case, the result remains a candidate until evidence gates support a bounded claim.

\begin{center}
\begin{tabular}{p{0.28\linewidth}p{0.26\linewidth}p{0.34\linewidth}}
\toprule
Backend & Status & Evidence Gate \\
\midrule
Classical PEC cavity path & Bounded internal Alpha evidence path & Residual and analytical reference comparison \\
C++ product layer & Planned or early Alpha infrastructure path & Contract and replay checks \\
DGTD backend & Future research route & DG residual, flux, boundary and conservation checks \\
PINN or neural-operator backend & Future research route & Maxwell residual, boundary, divergence, energy and reference gates \\
\bottomrule
\end{tabular}
\end{center}

DGTD and physics-AI backend support are future research routes. They are not presented as implemented production features in this public paper.

\section{AI-Assisted Solvers Are Untrusted Until Proven}

A future PINN or neural-operator backend should be treated as an untrusted field generator until it passes evidence gates. The AI is not trusted because it is neural. It is trusted only when its output survives Maxwell residual checks, boundary-condition checks, finite-value checks, reference comparison where available, and claim-boundary gates.

This is an AI-assisted solver quarantine. The quarantine is not a marketing label; it is a claim-control rule. Candidate fields, modes, traces, or spectra from a neural backend remain candidates until the evidence record states what was checked and what remains closed.

\section{C++ Product Direction and Artifact Replay}

C++ is the future product-core direction for mature DAD FieldWorks kernels and evidence-contract infrastructure. The current public-safe C++ direction is conservative: artifact replay, command-line skeleton concepts, result packaging, evidence rows, gate state, failure state, and claim boundaries.

In public specification terms, this direction includes a future C++ core interface and a future C++ implementation path for mature research kernels once evidence contracts, tests, and claim boundaries are strong enough to justify migration.

The public claim is architectural and replay-oriented. It does not claim a production CEM kernel. It does not claim arbitrary-geometry support. It does not claim an externally validated product. Private implementation source is not included in the public showcase repository.

\section{Public Showcase Visualizations}

The public website features selected public-safe visual diagnostics:

\begin{itemize}
  \item a main solver diagnostic hero animation,
  \item FDTD resonator ringdown with FFT spectrum,
  \item PEC cavity eigenmode field slice,
  \item Microwave Cavity Eigenmode Birth,
  \item Native Yee curl incidence microprototype.
\end{itemize}

The visuals are communication assets for diagnostic and architectural ideas. They do not broaden the claim boundary. Retained assets that are not featured on the homepage are not part of the current public visual story.

\section{Limitations and Claim Boundaries}

The public claim boundary is deliberately narrow:

\begin{itemize}
  \item no external validation claim,
  \item no production readiness claim,
  \item no commercial solver equivalence claim,
  \item no full-wave commercial solver claim,
  \item no arbitrary geometry support claim,
  \item no CPML support claim,
  \item no qubit simulation claim,
  \item no Josephson junction model claim,
  \item no processor-level quantum-hardware validation claim,
  \item no complete quantum hardware solver claim,
  \item no production C++ electromagnetic solver claim,
  \item no production AI solver claim,
  \item no implemented DGTD backend claim,
  \item no graphical application claim,
  \item no measurement validation claim.
\end{itemize}

These limitations are not cosmetic disclaimers. They are part of the evidence-contract model. A claim remains closed until the corresponding evidence path exists and passes.

\section{Roadmap}

The roadmap is evidence-led rather than feature-led.

\subsection{Near Term}

\begin{itemize}
  \item bounded evidence-contract cleanup,
  \item PEC cavity residual, mass, and divergence closure-readiness,
  \item native Yee curl incidence closure,
  \item public evidence documentation.
\end{itemize}

\subsection{Mid Term}

\begin{itemize}
  \item native Yee curl-curl planning,
  \item residual and mass policy resolution,
  \item C++ artifact replay and command-line contract layer,
  \item RF resonator diagnostics,
  \item classical solver evidence path hardening.
\end{itemize}

\subsection{Long Term}

\begin{itemize}
  \item DGTD backend route,
  \item physics-AI backend route,
  \item neural-operator and PINN quarantine under evidence contracts,
  \item external comparison pipelines when ready,
  \item graphical application only after the core evidence workflow matures.
\end{itemize}

No roadmap item should be read as already complete unless explicitly described as current public state.

\section{Conclusion}

DAD FieldWorks is being shaped as an evidence-contract framework for trustworthy computational electromagnetics. Its purpose is not merely to produce field plots, spectra, or eigenvalue candidates, but to bind such artifacts to residuals, comparisons, gates, decision classes, failure classes, risk records, and claim boundaries.

The current public state is useful but intentionally bounded. The public material presents a private research and development project through public architecture notes, public whitepaper material, and public-safe diagnostic visualizations. It does not publish private solver source code. It does not claim external validation. It does not claim production readiness. It does not claim commercial solver equivalence.

\section{References}

\begin{thebibliography}{9}

\bibitem{yee1966}
Kane S. Yee,
``Numerical solution of initial boundary value problems involving Maxwell's equations in isotropic media,''
\emph{IEEE Transactions on Antennas and Propagation}, 1966.

\bibitem{balanis}
Constantine A. Balanis,
\emph{Advanced Engineering Electromagnetics},
second edition, Wiley, 2012.

\bibitem{vvprinciple}
General verification and validation practice for computational models:
maintain separation between code verification, numerical diagnostics, reference comparison, uncertainty, and claim authorization.

\end{thebibliography}

\end{document}
